\section{From the spectral gap \texorpdfstring{\(4\to2\to6\to54\)}{4 to 2 to 6 to 54}: local coupling, modular compression, and two-dimensional integration}
\label{sec:sucesor83-salto-espectral-4-2-6-54}

For two consecutive representatives \((k,k+1)\), consider
\[
B_k=
\begin{pmatrix}
k^2&k(k+1)\\
k(k+1)&(k+1)^2
\end{pmatrix}\pmod 9.
\]

\begin{theorem}[Uniqueness of the Adjacent Diagonal Block]
\label{thm:sucesor83-app-bloque-central}
There is a unique block \(B_k\) whose diagonal entries coincide. It is
determined by \(k=4\) and is
\[
B_c=\begin{pmatrix}7&2\\2&7\end{pmatrix}.
\]
\end{theorem}

\begin{proof}
The equality of the diagonals is equivalent to
\[
k^2\equiv(k+1)^2\pmod 9,
\]
that is, \(2k+1\equiv0\pmod 9\). Since \(2\) is a unit of
\(\mathbb Z/9\mathbb Z\), the congruence has the unique solution
\(k\equiv4\pmod 9\).
\end{proof}

Therefore,
\[
B_c=7I+2R,
\qquad
R=\begin{pmatrix}0&1\\1&0\end{pmatrix},
\qquad R^2=I.
\]
With \(P_\pm=(I\pm R)/2\),
\[
B_c=9P_++5P_-.
\]

\begin{theorem}[Hierarchy of the separation between sum and product]
\label{thm:sucesor83-app-jerarquia-4-2-6-54}
The central block has local spectral jump \(9-5=4\). Its off-diagonal
coefficient is \(2\). Compression by the \(3\) classes modulo
\(3\) produces the aggregate difference \(51-45=6\), and expansion
over the \(9\) positions produces the global defect
\[
9\cdot6=459-405=54.
\]
The integers \(4,2,6,54\) pertain, respectively, to the spectral gap,
the local coupling, modular compression, and two-dimensional
integration.
\end{theorem}

\begin{proof}
The decomposition \(B_c=7I+2R=9P_++5P_-\) gives the jump and the local
coefficient. The aggregated matrices \(M_\Pi,M_\Sigma\) satisfy
\(\sum M_\Pi=51\) and \(\sum M_\Sigma=45\). Each aggregated entry represents
\(9\) positions, so unfolding multiplies the difference by
\(9\) and recovers the totals \(459\) and \(405\).
\end{proof}

Their ordered concatenations produce the blocks \(72\) and \(27\):
\[
72+27=99,
\qquad
72-27=45=\det B_c.
\]
The concatenation of the diagonal and antidiagonal entries produces,
respectively, \(77\) and \(22\). Therefore,
\[
72+27=77+22=99,
\qquad
77-22=55=54+1.
\]
These equalities express two exact partitions of the same block; they do not
by themselves identify an external constant.
